% page 195 du fac-simile (image p-194 du scan)
% bloc 165.1 x 224.7 mm ; marge gauche 36.9 mm
\pgc{15.240mm}{0.000mm}{
\l{\cel{56}187}
\l{}
\l{\sou{galvanizar}. (\sou{trans}.)Igar (muskuli inerta) kontraktar per la elek-}
\l{\cel{3}tro quan produktas en li la kontakto di du korpi heterogena.-}
\l{\cel{3}DEFIRS.}
\l{}
\l{\sou{galvanometro}. (\sou{fiziko)}. Instrumento qua uzesas por mezurar la}
\l{\cel{3}intenseso-grado di la fluo galvanala. - DEFIRS.}
\l{}
\l{\sou{galvanoplastar}. (\sou{trans}.)Aplikar strato de oro, arjento, zinko,}
\l{\cel{3}e c., per la efiko di baterio galvanala. - DEFIRS.}
\l{}
\l{\sou{galvanotropismo}. (\sou{biol.}) Efiko nemediata, da la elektro (fluo}
\l{\cel{3}kontinua), a la protoplasmo di ula celuli e, partikulare, di l}
\l{\cel{3}protozoi. - DEFIRS.}
\l{}
\l{\sou{gamo}. (\sou{muziko}).Serio de la sep noti di la skalo muzikala, dispo-}
\l{\cel{3}zita segun lia sucedo naturala, acensanta o decensanta, per}
\l{\cel{3}toni e mi-toni, intervale di un oktavo. - EFIRS.}
\l{}
\l{\sou{gambo}. I. (\sou{anat.}) Parto di la membro infra di la homo, qua pro-}
\l{\cel{3}longas la kruro, de la genuo til la pedo. - La tota membro. -}
\l{\cel{3}II. Che ula quadripedi : singla de la quar membri qui subtenas}
\l{\cel{3}la korpo e finas per hufi. - III. Membro di animalo quan lu}
\l{\cel{3}uzas por marchar, sizar, kaptar. - FI.}
\l{}
\l{\sou{gambitar}. (\sou{trans., ulu)}. Pasar sua gambo cirkum olta di persono,}
\l{\cel{3}por renversar ica. - FI.}
\l{}
\l{\sou{gambolar}. (\sou{netrans.}) Saltar agitante sua gambi. - EF.}
\l{}
\l{\sou{gamelo}. I. Granda skudelo en qua plura navani infra-grada, plura}
\l{\cel{3}soldati, kunmanjas. - II. Lada vazo en qua singla soldato re-}
\l{\cel{3}cevas sua porciono. - III. (\sou{milito-navo}). Gasto-tablo komuna}
\l{\cel{3}por la oficiri, la dicipuli, la kirurgiisti. - FIS.}
\l{}
\l{\sou{gameto}. (\sou{biol.})\cel{2}Celulo reproduktera, maskula o femina.}
\l{}
\l{\sou{gametofito}. (\sou{biol}.) Gameto di la planti, en la periodo en qua}
\l{\cel{3}la du pronuklei maskula e femina ne ja esas aglutinita (sporo-}
\l{\cel{3}fito) o kunfuzita (maxio).}
\l{}
\l{\sou{ganar}. (\sou{trans.}) Aquirar (profito-quanto). - EFIS.}
\l{}
\l{\sou{gango.}\cel{2}I. (\sou{teknol.}) Materio stranjera solida en qua kontenesas}
\l{\cel{3}erco o lapido. - II. (\sou{anat}.) Substanco amorfa qua envelopas}
\l{\cel{3}elemento anatomiala. - DEFIS.}
\l{}
\l{gangliono. (\sou{anat.}) Organo globatra ek tubero de fibri nervala}
\l{\cel{3}o de faskuli limfala. - DEFIRS.}
\l{}
\l{gangrenar. (\sou{netrans.)}\cel{1}(\sou{aludante la tisui animalala)}.Dissolvesar}
}
